%=========== ΛΥΣΗ - 116ο ΘΕΜΑ Β ===========

\begin{thema}{Β}
Στο σχήμα δίνονται οι γραφικές παραστάσεις των $f$ και $f^{-1}$, οι οποίες είναι συμμετρικές ως προς τη διχοτόμο $y=x$.

\begin{erwthma}
\item Από το σχήμα έχουμε
\[
f(2)=1{,}2.
\]
Άρα
\[
f^{-1}(1{,}2)=2.
\]
Για την παράγωγο της αντίστροφης ισχύει
\[
(f^{-1})'(y_0)=\frac{1}{f'(f^{-1}(y_0))}.
\]
Με $y_0=1{,}2$ παίρνουμε
\[
(f^{-1})'(1{,}2)=\frac{1}{f'(f^{-1}(1{,}2))}
=\frac{1}{f'(2)}
=\frac{1}{\frac{3}{5}}
=\frac{5}{3}.
\]

Η εφαπτομένη της $C_{f^{-1}}$ στο σημείο $A'(1{,}2,2)$ έχει εξίσωση
\[
y-2=\frac{5}{3}(x-1{,}2).
\]
Επειδή $1{,}2=\dfrac{6}{5}$, έχουμε
\[
y-2=\frac{5}{3}\left(x-\frac{6}{5}\right)
\Rightarrow y-2=\frac{5}{3}x-2
\Rightarrow y=\frac{5}{3}x.
\]

\item Από το σχήμα η $C_f$ βρίσκεται κάτω από τη διχοτόμο $y=x$ για κάθε $x\in(0,5)$ και τέμνει τη διχοτόμο μόνο στα άκρα $(0,0)$ και $(5,5)$. Άρα
\[
f(x)\leq x,\quad x\in[0,5],
\]
με ισότητα μόνο για $x=0$ και $x=5$.

Η $C_{f^{-1}}$ είναι συμμετρική της $C_f$ ως προς τη διχοτόμο $y=x$. Επομένως βρίσκεται πάνω από τη διχοτόμο για κάθε $x\in(0,5)$ και τέμνει τη διχοτόμο μόνο στα άκρα. Άρα
\[
f^{-1}(x)\geq x,\quad x\in[0,5],
\]
με ισότητα μόνο για $x=0$ και $x=5$.

\item Το ολοκλήρωμα $\dint_0^5 f(x)\d x$ εκφράζει το εμβαδόν κάτω από την $C_f$ στο $[0,5]$. Αντίστοιχα, το $\dint_0^5 f^{-1}(x)\d x$ εκφράζει το εμβαδόν κάτω από την $C_{f^{-1}}$ στο $[0,5]$.

Επειδή οι γραφικές παραστάσεις είναι συμμετρικές ως προς τη διχοτόμο $y=x$, τα δύο αυτά εμβαδά συμπληρώνουν το τετράγωνο πλευράς $5$ με κορυφές $(0,0)$, $(5,0)$, $(5,5)$, $(0,5)$. Άρα
\[
\dint_0^5 f(x)\d x+\dint_0^5 f^{-1}(x)\d x
=5\cdot 5=25.
\]

\item Από το Β2, στο $(0,5)$ ισχύει
\[
f(x)<f^{-1}(x).
\]
Άρα το εμβαδόν του χωρίου μεταξύ των $C_f$ και $C_{f^{-1}}$ είναι
\[
E=\dint_0^5\left(f^{-1}(x)-f(x)\right)\d x.
\]
Θέτουμε
\[
I=\dint_0^5 f(x)\d x.
\]
Από το προηγούμενο ερώτημα έχουμε
\[
\dint_0^5 f^{-1}(x)\d x=25-I.
\]
Επομένως
\[
E=(25-I)-I=25-2I.
\]
Άρα
\[
2I=25-E
\]
και τελικά
\[
\dint_0^5 f(x)\d x=\frac{25-E}{2}.
\]
\end{erwthma}
\end{thema}

%=========== ΤΕΛΟΣ ΛΥΣΗΣ - 116ο ΘΕΜΑ Β ===========
